% System design aligned with local first-release scope v0.2, 29 Oct 2026.
\chapter{SOLUTION DESIGN AND IMPLEMENTATION}\label{chap:design}

This chapter describes the target design confirmed on 29 October 2026. Baseline API/frontend code
exists, but the workspace, durable jobs and import extensions below are planned. A design element is
not a completed implementation or a passing test. The design follows the requirements in
Chapter~\ref{chap:problem}; existing implementation evidence is distinguished from planned work.

\section{System Overview}
\subsection{Product Context}
The research pipeline prepares data, trains and evaluates models offline. One local web application
loads a frozen model and reference snapshots; it supports researchers choosing substitutions for
follow-up experiments. Runtime input and project observations stay on the local computer.
\subsection{Functional Overview}
The first release covers UC-01 to UC-10 of Chapter~\ref{chap:problem}: prediction, ranking, export,
protein/tile lookup, experimental comparison, ClinVar/reference annotations, evidence pages,
peptide/protein jobs, experimental import and persistent project history. There is one product
configuration with local operation; public hosting and public sharing are excluded.

\section{System Architecture}
\subsection{Architectural Components}
Figure~\ref{fig:local-architecture} shows the target components. All application listeners use loopback;
reference downloads and training are preparation activities outside the runtime workflow.
\begin{figure}[htbp]
  \centering
  \begin{tikzpicture}[box/.style={draw,rounded corners,align=center,text width=3.25cm,
      minimum height=1.1cm,font=\small},flow/.style={-{Latex[length=2mm]}}]
    \node[box] (ui) {Local browser\\React + TypeScript};
    \node[box,right=0.7cm of ui] (api) {Local API\\FastAPI + inference};
    \node[box,below=0.8cm of api] (db) {Local PostgreSQL\\Reference + workspace};
    \node[box,left=0.7cm of db] (worker) {Local worker\\Shared inference package};
    \node[box,below=0.8cm of db] (files) {Local files\\Imports + frozen artefacts};
    \draw[flow] (ui) -- node[above,font=\footnotesize]{REST} (api);
    \draw[flow] (api) -- (db);
    \draw[flow] (worker) -- (db);
    \draw[flow] (worker.south) |- (files.west);
    \draw[flow] (api.east) -- ++(0.6,0) |- (files.east);
  \end{tikzpicture}
  \caption{Planned components of the single local application; arrows denote data interactions.}
  \label{fig:local-architecture}
\end{figure}

The React client owns input, matrix/ranking interaction and project/job/import screens. FastAPI
validates input, applies domain rules, serves reference annotations and stores snapshots. A local
worker uses the same inference package for jobs; PostgreSQL stores job state and project-owned
records. Generated local file paths hold committed imports and versioned artefacts. Reference ETL
never writes user observations into research data.

\subsection{Interaction Between Components}
A single-tile preview is validated and scored without persistence. Explicit Save creates an immutable
analysis in a selected project. Batch submission pins model, calibrator, gate and source versions,
returns a job ID, then the worker claims items transactionally. The client polls progress and opens
per-tile results. Expired worker leases are recovered after restart; completed items are not duplicated.
Cancellation and project deletion prevent late writes.

An import first produces a validation preview. Only a valid, explicitly committed dataset becomes
available for comparison. A comparison selects a saved analysis and dataset revision, checks assay
and mutation alignment, and stores its own result/provenance. Reading history uses saved snapshots,
not a new inference call or the latest ClinVar release.

\subsection{Design Rationale}
React supports the interactive matrix, filtered tables and long-running-task states. FastAPI keeps the
model computation in Python and provides an explicit REST contract. A local worker separates job
execution from interactive requests without introducing a cloud service. Persistent projects require
a database and snapshot lifecycle beyond the baseline transient preview. The trade-offs are local
resource limits, file/database consistency and additional recovery/deletion tests. SQLite remains
useful for isolated tests, while local PostgreSQL integration is an acceptance task. The OpenAPI
contract marks new operations as planned; generating a schema does not implement them.

\section{Data Design}
\subsection{Data Sources and Preprocessing}
Reference snapshots have version/date, checksum and terms of use. Research-only sources and
project-imported observations are kept distinct from served reference data (Table~\ref{tab:sources}).
\begin{table}[htbp]
  \centering
  \caption{Data roles in research and the first local release}\label{tab:sources}
  \small
  \begin{tabularx}{\textwidth}{|P{4.0cm}|X|}
    \hline \textbf{Source} & \textbf{Role and boundary} \\ \hline
    Zhang Mendeley v2 \citep{zhang2023data} & Mutation labels and reference observations; training use labelled. \\ \hline
    GEO GSE240610 \citep{geo240610} & GPS screen context; preserve estimator and source provenance. \\ \hline
    UniProtKB/Swiss-Prot & Pinned human reference sequences, accessions and coordinate/WT checks. \\ \hline
    ClinVar & First-release missense annotation; source/review/conflict state, attribution and diagnostic disclaimer. \\ \hline
    Project CSV/TSV & User observations for assay-aware comparison only; preserve source, units and replicates. \\ \hline
    Voutsinos/PAP \citep{voutsinos2026} & Research baselines and transfer-learning studies; not a replacement for mutation labels. \\ \hline
    Larsen \citep{larsen2025} and other restricted sources & Research-only consideration subject to rights; not automatically imported into the service. \\ \hline
  \end{tabularx}
\end{table}

\subsection{Data Model}
Reference entities are source, protein, tile, measured cell, ClinVar variant and model version.
The planned workspace adds project, analysis job, job item, analysis run/cell, imported dataset/observation
and comparison. Legacy expiring share tables are not reused as persistent project history.

\subsection{Data Dictionary}
\begin{table}[htbp]
  \centering
  \caption{Key workspace fields and invariants}\label{tab:workspace-fields}
  \small
  \begin{tabularx}{\textwidth}{|P{3.5cm}|X|}
    \hline \textbf{Entity} & \textbf{Required information} \\ \hline
    Project & UUID, name, description, creation/update time; one local workspace. \\ \hline
    Job and item & Mode/input IDs, pinned configuration, states/counts, item errors, coverage and worker lease; unique completed result per item. \\ \hline
    Analysis run/cell & Input and coordinate context; model/gate hashes; scores, rank, warnings, annotation/source snapshot and parent rerun ID. \\ \hline
    Dataset/observation & File hash/revision, assay/endpoint/units/source, WT sequence, position, WT/replacement, value/missing reason and replicate. \\ \hline
    Comparison & Run/dataset IDs, explicit replicate/aggregation, compatibility decision, alignment counts and metric/result snapshot. \\ \hline
  \end{tabularx}
\end{table}

\subsection{Relationships Between Entities}
A project owns jobs, runs and datasets. A job contains input items; each successful tile item produces
one immutable run with 532 cells. Dataset observations retain unique mutation/replicate keys.
A comparison joins one run and one committed dataset from the same project. Source/model artefacts
are versioned independently; changing active reference data must not change an old analysis.
Deletion shows dependencies, stops related work and removes owned files/results with retryable cleanup;
shared references and weights remain intact. Committed history has no automatic 30-day expiry.

\section{Method and Model Design}
Input remains the reference peptide, position and replacement; target is $\Delta$PSI. The existing
service baseline is the frozen ridge/Platt model from Section~\ref{sec:res-calibration}. This scope
update does not choose a replacement model or change the research comparisons with PAP/Degformer.
Raw effect and calibrated $P(\Delta\mathrm{PSI}\geq1)$ remain separate fields.

\paragraph{Domain integration.}
The effect model's intrinsic-WT training threshold must not be applied directly to the GPS screen
estimator. The R3 research provides a two-tier gate with an intermediate borderline state; integrating
its frozen rules and provenance into the baseline remains pending. Every tile has domain evidence;
known out-of-domain scoring requires an explicit warned override. Whole-protein mode is a collection
of peptide-context predictions, not a validated intact-protein model.

\section{Functional Design}
\subsection{Batch and Protein Processing}
Peptide mode accepts a bounded list or multi-record FASTA of 28-residue inputs. Protein mode accepts
one reference accession or one FASTA record. Tile $n$ starts at $(n-1)\cdot23+1$ and ends 27 residues
later; only full windows are generated. Report uncovered intervals, invalid windows and domain
refusals. Overlapping protein substitutions retain their tile contexts without averaging.
Job states are pending, running, completed, partial, failed or cancelled; per-item states and counts
explain partial results. Browser tables are paginated; exports include every outcome.

\subsection{Import and Comparison}
CSV/TSV rows contain peptide ID, WT sequence, 1-based position, WT/replacement, observed value and
replicate. Optional accession/release/start coordinates are checked together. Metadata identify assay,
endpoint, units, source and reference. Duplicate/conflicting mutation-replicate keys, nonfinite values
and WT mismatches block commit. Empty observations retain a missing reason.

Comparison requires exact mutation alignment and explicit replicate selection or aggregation.
Comparable GPS $\Delta$PSI supports residuals (prediction minus observation), MAE and Spearman with
sample counts; insufficient/constant pairs have no reported correlation. Different assay endpoints
remain separate without pooled metrics. Imported observations do not automatically constitute
independent validation or become training labels.

\subsection{Annotations and Project Screens}
The first release shows ClinVar alongside approved UniProt/GPS context. Show classification,
review/conflict state, source release and WT-coordinate agreement; distinguish no match from
unavailable/ambiguous mapping. Attribution and diagnostic disclaimer travel with saved/exported
annotations. Custom unmapped sequences have explicit annotation-unavailable states.
The project workspace exposes history, jobs and datasets; users can reopen snapshots, create reruns
as new entries and confirm deletions. New screens and API routes remain pending implementation.

\section{Local Execution View}
One computer hosts the browser, API, worker, PostgreSQL and workspace files. Application listeners
bind to loopback; the frontend origin is explicitly allowed. Docker may package the same components,
but no remote host, public URL, TLS deployment or separate product edition is required.
Appendix~\ref{app:installation} separates baseline startup from the pending first-release workflow.
